(1) Over $\Mo$ the form is nondegenerate; diagonalising at $x$ gives $A(x) \cong \Cl(p', q')$ for the local signature $(p', q')$, which is $(p,q)$ on the physical side and $(p{-}1, q{+}1)$ across the wall (Remark \ref{rem:coords}). The named isomorphisms are Fact \claimref{sheaf.imported-classification-values} (the coincidence of the two is the fusion phenomenon of \cite[Ex.~5.10]{ButlerT1}, both being ascents of $\Cl(0,3) \cong \HH \oplus \HH$; it plays no role here).

(2) Since $k$ spans the radical, $G_x(k, \alpha) = 0$ for \emph{every} $\alpha \in T^*_xM$, so $k$ anticommutes with every covector: $k\alpha + \alpha k = 2\,G_x(k,\alpha) = 0$ (\claimref{std-algebra.vector-squares}). Hence for every $a \in A(x)$,
\begin{equation}\label{eq:kcommutation}
k\,a \;=\; \hat{a}\,k ,
\end{equation}
with $\hat{\cdot}$ the grade involution (the commutation rule \claimref{primer.commutation-rule}, for a $k$ orthogonal to everything): on a product of generators, $k\,(\alpha_1 \cdots \alpha_j) = (-1)^j (\alpha_1 \cdots \alpha_j)\, k = \widehat{\alpha_1 \cdots \alpha_j}\; k$ by $j$ anticommutations, and both sides of \eqref{eq:kcommutation} are linear in $a$ (nondegeneracy plays no role). The fibre statements of (2) are the case $r = 1$ of Theorem \claimref{std-algebra.radical-ideal-and-quotient}: its part (i) gives $J(x) = Ak = kA$ of dimension $2^{n} - 2^{n-1}$, its part (ii) gives $J(x)^{2} = 0$, and its part (iii) gives the quotient, the projection onto $T^*_xM/\shL(x)$ extending to a surjective homomorphism with kernel exactly $J(x)$. In the notation used below: $A k = k A$, and this common set is a two-sided ideal containing $k$; conversely any two-sided ideal containing $k$ contains $Ak$, so $J(x) = Ak$. It squares to zero:
\[
J(x)^2 \;=\; A\,k\,A\,k \;=\; A\,\hat{A}\,k^2 \;=\; 0 ,
\]
by \eqref{eq:kcommutation} and $k^2 = G_x(k,k) = 0$. For the dimension, choose a complement $V_1 \subset T^*_xM$ of $\shL(x)$ and take a basis of $V_1$ together with $k$; ordering $k$ last, the blade basis of Lemma \claimref{sheaf.blade-basis-degenerate-forms-included} splits $A(x) = \Cl'(V_1) \oplus \Cl'(V_1)\,k$, where $\Cl'(V_1)$ denotes the span of the blades in the $V_1$-letters (a subalgebra, since $V_1$ is closed under the rewriting against itself). Then $Ak = \Cl'(V_1)\,k$, of dimension $2^{n-1}$: indeed $(\,b\,k\,)k = \pm b k^2 = 0$ for $b \in \Cl'(V_1)$, so the $\Cl'(V_1)k$-part of $A$ contributes nothing new, and left multiplication by elements of $\Cl'(V_1)$ preserves $\Cl'(V_1)k$.

For the quotient: the projection $\pi : T^*_xM \to T^*_xM/\shL(x)$ intertwines $G_x$ with a well-defined form $\overline{G}_x$ ($G_x(\alpha + k, \beta + k') = G_x(\alpha, \beta)$ for $k, k'$ radical), which is nondegenerate of rank $n - 1$: any $\alpha$ with $\overline{G}_x(\bar\alpha, \cdot) = 0$ has $G_x(\alpha, \cdot) = 0$, hence $\alpha \in \shL(x)$, hence $\bar\alpha = 0$. Its signature is $(p{-}1, q)$: on the physical side the full form diagonalises with entries $(+^{p}, -^{q})$ degenerating so that one $+$ goes to zero on $\Sig$ (the dying eigenvalue is positive there, Remark \ref{rem:coords}), leaving $(p{-}1, q)$ on the quotient by continuity of the nondegenerate block. The composite $T^*_xM \xrightarrow{\pi} T^*_xM/\shL(x) \hookrightarrow \Cl(T^*_xM/\shL(x), \overline{G}_x)$ satisfies the Clifford relation for $G_x$, so by the universal property (\claimref{std-algebra.existence-uniqueness}) it extends to a surjective algebra map $\Pi : A(x) \to \Cl(\overline{G}_x)$ (surjective because the generators are in the image). Its kernel contains $k$, hence contains the two-sided ideal $J(x)$; and by dimensions the containment is equality:
\[
\dim A(x) - \dim J(x) \;=\; 2^n - 2^{n-1} \;=\; 2^{n-1} \;=\; \dim \Cl(\overline{G}_x),
\]
using Lemma \claimref{sheaf.blade-basis-degenerate-forms-included} on the quotient form. So $A(x)/J(x) \cong \Cl(\overline{G}_x)$, canonically ($\Pi$ used no choices).

(3) Working locally on $\Sig$, choose a spanning section $k$ of $\shL$ (a smooth line subbundle by \claimref{sheaf.tameness}) and a local frame of $T^*M|_\Sig$ completing it; the sections $\{e_I\, k : k \notin e_I\}$ ($2^{n-1}$ of them, the blades in the complementary letters times $k$) are smooth and pointwise a basis of $J(x)$ by (2), so the image of the multiplication map is a smooth subbundle of rank $2^{n-1}$, independent of the choices (the fibre $J(x)$ is choice-free). Its pushforward $\shJ$ is finitely presented: locally, with $\tau$ a defining function of $\Sig$, the presentation $\mathcal{O}^{2^{n-1}} \xrightarrow{\ \tau\ } \mathcal{O}^{2^{n-1}} \to \shJ \to 0$ exhibits it (sections of the bundle over $U \cap \Sig$ are smooth functions modulo those vanishing on $\Sig$, componentwise in the displayed frame). It is supported exactly on $\Sig$, since its stalks vanish off $\Sig$ and are nonzero on it, and it is an $\shA$-bimodule because each fibre is a two-sided ideal and the module operations are the restrictions of the smooth multiplication of $\shA$. For $\shK$: the map $\shA \to i_*(\shA|_\Sig / \shA|_\Sig \cdot \shL)$ (restrict to $\Sig$, then project) is a morphism of $\shA$-bimodules, since each fibre quotient is a quotient by a two-sided ideal; its kernel is therefore a sheaf of two-sided ideals of $\shA$, its sections over $U$ are exactly the sections of $\shA$ whose restriction to $U \cap \Sig$ lies in the subbundle, and the stated fibre and quotient identifications follow from (2).
